\begin{enumerate}
	\item On développe l'expression factorisée proposée :

	$(0,5x + 1,5)(-x+5) = 0,5x \times (-x) + 0,5x \times 5 + 1,5 \times (-x) + 1,5 \times 5 $

	On réduit les produits :

	$(0,5x + 1,5)(-x+5) = -0,5x^2 + 2,5x - 1,5x + 7,5 $

	On réduit les sommes :

	$(0,5x + 1,5)(-x+5) = -0,5x^2 + x + 7,5 $

	On obtient bien l'expression de $f(x)$ donnée dans le sujet.

	Pour tout $x$ appartenant à l'intervalle $[-4 ; 6]$,
	\[
	f(x)=(0,5x + 1,5)(-x+5) = -0,5x^2 +  x + 7,5.
	\]

	\item Le nombre dérivé en 3 est la pente de la tangente à la courbe au point d'abscisse 3.

	Par lecture graphique, \quad $f'(3) = \dfrac{4 - 6}{4 - 3} = \dfrac{-2}{1} = -2$.

	En donnant l'expression de la fonction dérivée : \quad $f'(x) = -0,5 \times (2x) + 1 = -x+1$.

	On détermine \quad $f'(3) = -3  + 1 = -2$.


	\item \begin{enumerate}
		\item Graphiquement, les points de la courbe d'ordonnée inférieure ou égal à 0 sont les points d'abscisse appartenant à $\left[-4 ; -3\right]\cup\left[5 ; 6\right]$.

		Algébriquement, on utilise la forme factorisée pour déterminer le signe.

		$\begin{aligned}[t]
			0,5x + 1,5 = 0
			&\iff 0,5x = -1,5 \\
			&\iff x = -3
		\end{aligned}$
		\quad et \quad
		$\begin{aligned}[t]
			0,5x + 1,5 > 0
			&\iff 0,5x > -1,5 \\
			&\iff x > -3
		\end{aligned}$

		$-x+5 = 0 \iff x = 5$\quad  et  \quad $\begin{aligned}[t]
			-x+5 > 0 &\iff 5 > x \\&\iff x < 5.
		\end{aligned}$

	Sur l'intervalle $[-3 ; 5]$,\quad $0,5x + 1,5 \geqslant 0$ \quad  et \quad $-x+5 \geqslant 0$

	donc le produit \quad $(0,5x + 1,5)(-x+5)$ \quad est positif.

	\textbf{L'affirmation~a est fausse.}

	\item Graphiquement, on voit que la fonction est décroissante sur l'intervalle $[1 ; 4]$, donc le nombre dérivé $f'(x)$ est négatif sur l'intervalle $[1 ; 4]$.

	On a vu que \quad $f'(3) = -2$ \quad donc $f'(x)$ n'est pas positif sur l'intervalle $[1 ; 4]$.

	\textbf{L'affirmation~b est fausse.}
	\end{enumerate}
\end{enumerate}
